\chapter{Metodologia eksperymentalna}

\section{Przygotowanie danych}

Przed właściwym procesem trenowania modeli dane zostały poddane wstępnej obróbce. W analizowanym zbiorze nie występowały braki danych, więc ich filtracja nie była konieczna. Ze względu na obecność wartości odstających zmiennej \texttt{playtime\_forever} zdecydowano się na ograniczenie jej wartości do $120000$ minut, co odpowiada 2000 godzinom, oraz w przybliżeniu $99{,}97$ percentylowi w zbiorze danych. Obserwacje powyżej tej granicy zostały zastąpione wartością 120000.  Założono, że wartości powyżej tej granicy nie niosą istotnej informacji o preferencjach graczy, a jedynie mogą zaburzyć proces trenowania modeli.

Usunięto również gry przypisane do mniej niż 5 użytkowników oraz użytkowników posiadających w swojej bibliotece mniej niż 5 gier. Miało to na celu zmiejszenia wpływu zimnego startu na wyniki modeli oraz możliwość późniejszego podziału danych na zbiory treningowy, walidacyjny i testowy, w sposób zapewniający, że każdy użytkownik pojawi się w każdym z tych zbiorów.

Zbiór po filtracji został podzielony na zbiory: treningowy, walidacyjny i testowy. Dla każdego użytkownika do zbioru testowego przypisano liczbę interakcji równą
$\left\lceil 0{,}1 \cdot n_u \right\rceil$,
gdzie $n_u$ oznacza liczbę interakcji danego użytkownika. Zastosowano zaokrąglenie w górę, aby zapewnić, że każdy użytkownik będzie miał co najmniej jedną interakcję w zbiorze testowym. Podział został przeprowadzony w sposób losowy ze względu na brak inforamcji o datach przypisania produktów do bibliotek użytkowników. Pozostały zbiór został analogicznie podzielony na zbiór treningowy i walidacyjny. W celu zachowania reprodukowalności wyników, na każdym etapie podziału zostało ustalone ziarno generatora liczb pseudolosowych.

\section{Porównywane modele}
    
Porównano sześć modeli rekomendacyjnych, w tym cztery modele implicit ALS z różnymi funkcjami wag pewności (confidence), model Bayesian Personalized Ranking (BPR) zawarte w bibliotece implicit \texttt{implicit} \cite{frederickson2022implicit} oraz model popularnościowy jako punkt odniesienia. Zastosowane funkcje wag pewności w modelach implicit ALS to:
\begin{itemize}
    \item $c_{ui} = 1 + \alpha r_{ui}$,
    \item $c_{ui} = 1 + \alpha \log(1 + r_{ui})$,
    \item $c_{ui} = 1 + \alpha \frac{\log(1 + r_{ui})}{1 + \log(1 + r_{ui})}$,
    \item $c_{ui} = 1 + \alpha (\log(1 + r_{ui}) + \beta)$,
\end{itemize}

gdzie $r_{ui}$ oznacza zmienną \texttt{playtime\_forever} dla użytkownika $u$ i gry $i$, a $\alpha$ jest domyślnym hiperparametrem w modelach implicit ALS. Dwie pierwsze funkcje wag pewności zostały zaproponowane w pracy \cite{hu2008implicit}, natomiast kolejne zostały utworzone na podstawie wyników poprzednich modeli i analizy rozkładu zmiennej \texttt{playtime\_forever} w zbiorze danych. W szczególności funkcja $c_{ui} = 1 + \alpha (\log(1 + r_{ui}) + \beta)$ została zaproponowana w celu zmniejszenia wpływu szumu wynikającego z dużej liczby obserwacji o niskiej wartości zmiennej \texttt{playtime\_forever}, które ze względu na własności logarytmu mają stosunkowo wysoką wartość w funkcjach wag pewności, a nie są realnym odzwierciedleniem preferencji użytkowników. Dodatkowo przesunięcie  $\beta$ pozwala nadać obserwacjom z \texttt{playtime\_forever = 0} wagę większą niż 1, co miało na celu wyrazić zainteresowanie użytkownika grą, wynikające z jej zakupu, mimo braku czasu spędzonego w grze. Aby uniknąć arbitralnego wyboru wartości $\beta$, w eksperymencie przeprowadzono strojenie tego hiperparametru.

\section{Strojenie hiperparametrów}



\section{Finalna ewaluacja testowa}

\section{Analiza stabilności, wrażliwości i testy statystyczne}

\section{Ograniczenia metodologii}